% Generated by records_reorg.py assemble (records-reorg/); edit the data and re-run, not this file.
\subsection{Data, simulation and event selection}\label{recsec:method-data}

% Compression rewrite (DECISION-20261119c; REWRITE-RULES-20261119.md rule 5): text the limitation-resolver lane
% moved verbatim out of the posted note, with one move-manifest row per sentence (LR-C ids in
% history-cuts/move-manifest-lr-20261117.tsv). References into the note resolve through main_note.aux (xr).
% [records reorg] former heading: \section{Method and statistics detail moved from the note}
% [records reorg, not printed] This section holds the procedure and supporting detail that the compressed note no longer prints, as the note printed it.

\paragraph{Per-event analysis trees.}
% [records reorg] former heading: \subsection{Per-event analysis trees (from Part 2, sec:evloop)}
% From sec_experiment.tex:34 (draft f6a4fc372 lineage); procedure a reader does not need to interpret a result
Unbinned unfolding needs per-event records rather than histograms. A dedicated event loop (App.~\ref{app:codebase}), built on MAT's systematic-universe tools, writes four per-event samples for each playlist: all true signal events (the efficiency denominator), reconstructed truth-matched simulated events, simulated background, and data. Entries carry truth and reconstructed observables ($\pT$, $\pPar$, and in extended runs $\Eavail$, $q_{3}$, $W$) and the MINERvA Tune~v1 central-value weight. The systematic-universe samples also carry every universe's weight and the shifted observables of each lateral band, a variation that shifts reconstructed kinematics rather than reweighting events (a vertical band; \S\ref{sec:systematics}). Two conditions make the truth sample the complete reference for the simulated signal:
% (\texttt{\_universes\_full})
\begin{itemize}
  \item A reconstructed signal event is filled only if its $(\mathrm{run},\mathrm{subrun},\mathrm{event})$ key is in the truth denominator. Truth events without a reconstructed partner enter as native OmniFold ``miss'' entries.
  \item The simulated signal sample is complete by construction at ME~FHC (completeness $c=1.000000$; \S\ref{sec:method-xsec}).
\end{itemize}
The loop reads one playlist-wide flux and calibration from the first event, so it is run separately for each playlist, and the twelve per-playlist outputs are combined. The corresponding baseline histogram event loop supplies the flux histogram, POT-weighted across playlists; the size of the playlist-to-playlist flux differences is stated in \S\ref{rlim:experiment}.
% (\texttt{TTree::SetMaxTreeSize}-aware)

\paragraph{Shapes of the reconstructed control distributions.}
% [records reorg] former heading: \subsection{Shapes of the reconstructed control distributions (from Part 2, sec:controlplots)}
% From sec_experiment.tex:49 (draft f6a4fc372 lineage); description of the control figures beyond their captions
The beam and acceptance set the shapes. The $\pPar$ peak at $4$--$6$~GeV/$c$ reflects the ${\langle E_\nu\rangle}\approx\SI{6}{GeV}$ NuMI medium-energy peak; spectrometer matching removes the lowest-$\pPar$ muons. The $\pT$ spectrum peaks near $0.7$~GeV/$c$ and falls by over an order of magnitude across the measured range. Low-recoil (quasi-elastic and $2p2h$) events dominate the sample. Available energy $\Eavail$ concentrates below ${\sim}\SI{0.5}{GeV}$, with a long inelastic tail in the wide upper bins; $q_{3}$ tracks this population's momentum transfer. Reconstructed $W$ rises through the smeared quasi-elastic and $\Delta(1232)$ region before the deep-inelastic continuum takes over above ${\sim}\SI{2}{GeV}$. Absolute background is at the few-percent level under the muon projections and concentrates at high recoil (highest $\Eavail$, $q_{3}$ and $W$ bins).

\paragraph{Detector layout and muon spectrometer.}
% [records reorg] former heading: \subsection{From Part 2, secexperiment.tex: the paragraph beginning ``MINERvAciteMINERvA:2013zvz was a fine-grained scintillator-tracker neutrino''}
% From sec_experiment.tex:5 @d55f7aad6; qualifier detail or procedure (DECISION-20261120a, second pass)
In the paragraph as the note printed it, this text followed: ``MINERvA~\cite{MINERvA:2013zvz} was a fine-grained scintillator-tracker neutrino experiment on Fermilab's NuMI beamline.''\par
Its central tracker comprises hexagonal planes of triangular polystyrene scintillator strips in three stereo orientations, surrounded by electromagnetic and hadronic calorimetry. The magnetized MINOS near detector, \SI{2}{m} downstream, serves as the muon spectrometer, measuring momentum and charge for MINOS-matched tracks.

\paragraph{Reweighting components of the reference simulation.}
% [records reorg] former heading: \subsection{From Part 2, secexperiment.tex: the paragraph beginning ``The reference simulation uses GENIE 2.12.6citeAndreopoulos:2009rq''}
% From sec_experiment.tex:10 @d55f7aad6; qualifier detail or procedure (DECISION-20261120a, second pass)
In the paragraph as the note printed it, this text followed: ``The reference simulation uses GENIE 2.12.6~\cite{Andreopoulos:2009rq} and MINERvA Tune v1, applied as per-event weights by the reweighting tools of the MINERvA Analysis Toolkit (MAT).''\par
These are the flux-and-central-value (flux-and-CV) weight, the GENIE model-parameter (``dial'') weights, the low-recoil 2p2h enhancement fit, and the non-resonant-pion reduction.

\paragraph{Background subtraction before unfolding.}
% [records reorg] former heading: \subsection{From Part 2, secexperiment.tex: the paragraph beginning ``Genuine backgrounds are selected events whose''}
% From sec_experiment.tex:26 @d55f7aad6; qualifier detail or procedure (DECISION-20261120a, second pass)
In the paragraph as the note printed it, this text followed: ``Out-of-phase-space signal, reconstructed inside the phase space with its true muon outside it, is routed into the background template with them.''\par
Both are POT-scaled and subtracted from the reconstructed-data target before unfolding (\S\ref{sec:method}). Their sum, $N_{\mathrm{bkg}}(b)$ in each reconstructed bin, enters the purity weight $w_{\mathrm{pur}}(b)$ of \S\ref{sec:method-algo}.

\paragraph{Composition of the background template.}
% [records reorg] former heading: \subsection{From Part 2, secexperiment.tex: the paragraph beginning ``The background template used for this''}
% From sec_experiment.tex:28 @d55f7aad6; qualifier detail or procedure (DECISION-20261120a, second pass)
% [lead-in, Option B] In the paragraph as the note printed it, this text followed: ``The background template used for this result is \SI{2.91}{\percent} of selected data events in the two-dimensional grid over the full sample (\SIrange{2.72}{3.23}{\percent} per playlist, \SI{2.74}{\percent} on playlist 1A).''\par
On playlist 1A only, about \SI{10}{\percent} (weighted) of the template is out-of-phase-space signal. These are $\nu_{\mu}$ charged-current events with a true vertex inside the fiducial volume. The rest is genuine background, mostly $\nu_{\mu}$ charged-current events with a true vertex outside it.

\paragraph{Data-to-simulation normalization at reconstructed level.}
% [records reorg] former heading: \subsection{From Part 2, secexperiment.tex: the paragraph beginning ``Selected data exceed the POT-scaled MINERvA''}
% From sec_experiment.tex:39 @d55f7aad6; qualifier detail or procedure (DECISION-20261120a, second pass)
Selected data exceed the POT-scaled MINERvA Tune~v1 prediction (signal $+$ background) by $\sim\SI{12}{\percent}$ in overall normalization.
% [records reorg] former heading: \subsection{From Part 2, secexperiment.tex: the paragraph beginning ``Selected data exceed the POT-scaled MINERvA''}
% From sec_experiment.tex:39 @d55f7aad6; qualifier detail or procedure (DECISION-20261120a, second pass)
% [lead-in, Option B] In the paragraph as the note printed it, this text followed: ``The data/MC ratio rises with $\pT$, from $0.88$ in the lowest bin to $\sim1.2$ above \SI{1}{GeV/c}.''\par
Excess and trend together reflect the truth-level normalization deficit of \S\ref{sec:results-model} and \S\ref{sec:3d-models}. Beyond the $\pT$ trend noted above, the ratio panels also show a mild enhancement in the lowest $\Eavail$ bins, reflecting the low-recoil excess examined at truth level in \S\ref{sec:3d-2p2h}.

\paragraph{Status of the PET method.}
% [records reorg] former heading: \subsubsection{From Part 2, the selection-and-classifier table caption (row n06)}
% From selection_features.tex:6 @e89694209; register row, check stage 2026-11-21
% absence-ok: verbatim move of a register sentence; the note's PET pointer is sec6_conclusions.tex (point-cloud event representations)
% [lead-in, Option B] In the paragraph as the note printed it, this text followed: ``Selection and classifier information are separate choices. Reporting bins are not classifier inputs.''\par
PET remains diagnostic method development; its design-study recovery does not validate real-data background subtraction.

\subsection{Observable definitions}\label{recsec:method-observables}

\paragraph{Truth-level W residual.}
% [records reorg] former heading: \subsection{Truth-level W residual (from Part 2, the five observables)}
% From sec_observables.tex:40 (draft f6a4fc372 lineage); detail a reader does not need to interpret a result
At truth level, $q_{3}$ is computed from the true neutrino and muon four-vectors, but $\theta_{\mu}$, and with it $\pT$ and $\pPar$, is measured from a fixed beam axis tilted by \SI{58.87}{mrad} from the detector axis rather than from the true neutrino direction.  Truth $W$ is therefore fixed by $(\pT,\pPar,q_{3})$ only up to that event-by-event angle, of about \SI{1}{mrad}, and the choice of $M$: the residual is within \SI{\pm2}{MeV} for 68\% of events and within $-49$ to $+67$~MeV for 99.9\%; \SI{0.25}{\percent} of events (weighted) change their coarse $W$ bin, and \SI{0.086}{\percent} (weighted) have no real solution. % D1 (crossdoc audit 2026-11-15): weighted, probe_omni.json:10 q3sq_minus_pt2_negative_wfrac 0.0008646 (the 0.15% was the unweighted 49754/32846302, :9 and :8); the Supplement prints m4w-noreal-pct, 0.086
% Methods rewrite (M4, the truth-W residual): M4-W-RECORD.md :6 (theta_true from a fixed beam axis tilted -58.87 mrad
% about x, MAT TruthFunctions.h:21-27, PlotUtilsPhysicalConstants.h:44), :7 and :74-81 (W_true - W_pred with M_bar: 68%
% within about +-2 MeV, 99.9% within -49/+67 MeV), :85-90 (0.25% of events, weighted, change coarse W bin; 0.15% have
% no real solution), :63 (the tilt is the only other cause), :55-59 (median 1.03 mrad between the true neutrino and that
% axis); q3 from the true four-vectors, README-gen5d-build.md:50-58. REVIEW-1 finding 18: M4-W-RECORD.md:7, :72.

\paragraph{Energy transfer $q_{0}$.}
% [records reorg] former heading: \subsection{From Part 2, secobservables.tex: the paragraph beginning ``Tablereftab:observables-impl gives the implemented truth and''}
% From sec_observables.tex:6 @d55f7aad6; qualifier detail or procedure (DECISION-20261120a, second pass)
% [lead-in, Option B] In the paragraph as the note printed it, this text followed: ``$E_{\mu}$, $p_{\mu}$, and $\theta_{\mu}$ are the muon energy, momentum, and angle.''\par
$q_{0}$ is the calorimetric estimate of the energy transfer.

\paragraph{Classifier inputs and reported observables.}
% [records reorg] former heading: \subsection{From Part 2, secobservables.tex: the paragraph beginning ``The five reconstructed inputs therefore share''}
% From sec_observables.tex:41 @d55f7aad6; qualifier detail or procedure (DECISION-20261120a, second pass)
% [lead-in, Option B] In the paragraph as the note printed it, this text followed: ``The five reconstructed inputs therefore share information and do not, by themselves, provide five independent constraints: the result is five simultaneous observables, not five independent measurements.''\par
The five classifier inputs, the lower-dimensional support they populate, and the five reported observables are distinct. Only reconstructed $W$ is fixed by the other inputs; reconstructed $\Eavail$ comes from a separate calorimetric estimator and is not a function of them. The nearly degenerate correlation directions of the joint-test metric (\S\ref{sec:joint5d-design}) are consistent with this redundancy, but that origin has not been shown, and the information added by $W$ has not been measured.
% carried with the text from sec_observables.tex:42 @d55f7aad6
% absence-ok: "has not been shown" and "has not been measured" are the base's own limitations of the five inputs, reworded in r22 Step 2

\paragraph{Effect of the sample-definition convention on cells.}
% [records reorg] former heading: \subsection{From Part 2, secobservables.tex: the paragraph beginning ``The analysis excludes events with negative''}
% From sec_observables.tex:44 @d55f7aad6; qualifier detail or procedure (DECISION-20261120a, second pass)
% [lead-in, Option B] In the paragraph as the note printed it, this text followed: ``The generators measure $\pT$ and $\pPar$ about their own neutrino axis; their $\pT$ differs from the analysis truth $\pT$ by \SI{\pm4}{MeV} for 68\% of events.''\par
Over the 109 joint-test cells, this convention changes the analysis simulation's cell contents by factors between 0.9965 and 1.0035 (median 1.000008), at the level of its Monte Carlo precision, and shifts GENIE's central prediction by between $-0.47$ and $+0.53$ of each cell's null standard deviation.
% carried with the text from sec_observables.tex:45 @d55f7aad6
% Search: information added by W; W redundancy; near-degenerate directions; origin; joint-test metric; pT axis; neutrino axis; beam axis; MP-5; R_c (none for W's information or the origin: M2-ARCHIVE-CHECKS.md:7, r12/ledger.md:103-104 restate it; the pT-axis effect is printed from MP-5, state/measurement-paper/phase1/mp5/REPORT.md:66, :68, :100, job 60017887, with its transfer to the other generators UNRELIABLE, :87-90; Joseph, DECISION-20261027-joseph-absence-audit-items-3-4.md item 1) @ 7a6ae651

\paragraph{Charged-pion mass replacement in $E_{\mathrm{avail}}$.}
% [records reorg] former heading: \subsubsection{From Part 2, the charged-pion offset in the available energy}
% From sec_3d.tex:36 @e89694209; check-stage claims read
% Search: 135 MeV; 139.57; mass_pion; pion mass; OI-30; UD-002 (stands: DECISION-20261101-joseph-eavail-keep-135.md item 1); materiality; covariance projection (done for the sixteen-cell result: ASSESSMENT-20261031-eavail-charged-pion-mass-switch-pricing.md sec. 7.3, full-covariance norm 0.43 prior level / 0.42 first order; EAVAIL_DEFINITION.md sec. 3 as corrected by #637) @ c93a73bce
% [lead-in, Option B] In the paragraph as the note printed it, this text followed: ``The charged-pion term subtracts the inherited constant \SI{135}{MeV} instead of the charged-pion mass of about \SI{139.57}{MeV}.''\par
Relative to the charged-pion mass, it moves \fvval{eavbin-share}\,\% of simulated signal weight to other bins on the fine grid of seven true-$\Eavail$ bins and changes none of the seven bins by more than \fvval{eavbin-max}\,\%. This effect has not been assessed against the projected five-dimensional covariance of the \resFiveAxis{} from the \estPurity{} (estimator definitions, Table~\ref{tab:estimators-matrix}). We tested the mass replacement on the \resSixteen{} (Sec.~\ref{sec:fived16}) from the \estMultiD{}. The cross section shifts by at most \fvval{eav-cellmax}$\sigma$ in any cell. The coherent shift has norm \fvval{eav-coherent}$\sigma$ under the reported covariance.

\paragraph{Generator neutrino axis for muon kinematics.}
% [records reorg] former heading: \subsubsection{From Part 2, the generators' muon axis}
% From sec_observables.tex:44 @e89694209; check-stage claims read
The generators measure $\pT$ and $\pPar$ about their own neutrino axis; their $\pT$ differs from the analysis truth $\pT$ by \SI{\pm4}{MeV} for 68\% of events. For GENIE with MEC, NuWro and GiBUU its effect is not reliably determined (Ref.~\cite{NoteRecords2026}, \recentry{mp5-truth-axis}).

\subsection{The OmniFold implementation and its training settings}\label{recsec:method-omnifold}

\paragraph{Iteration count and the Huang et al.\ recommendation.}
% [records reorg] former heading: \subsection{Iteration count and the Huang et al.\ recommendation (from Part 2, sec:method-config)}
% From sec_method.tex (dfc41cdf5) (draft f6a4fc372 lineage); rationale a reader does not need to interpret a result
Huang \emph{et al.}~\cite{Huang:2025omnifold} select 45 iterations for OmniFold-family methods (20 for their binned benchmark) using a binned $\chi^{2}$ against known mock-data truth, a criterion they note is infeasible on real data. We do not adopt that recommendation. Their mock data injects a large prior--data \emph{shape} discrepancy through an RPA-modified quasielastic model, requiring many iterations. The simulation prior differs from the measured distributions in shape as well as normalization: the data exceed it by about 12\%, and the data/MC ratio rises with $\pT$ from 0.88 to about 1.2 (\S\ref{sec:controlplots}). Cross-experiment production also typically uses ${\sim}5$ iterations~\cite{Canelli:2025guide}.

\paragraph{MINERvA driver and pipeline conventions.}
% [records reorg] former heading: \subsection{MINERvA driver and pipeline conventions (from Part 2, sec:method-driver)}
% From sec_method.tex (dfc41cdf5) (draft f6a4fc372 lineage); procedure a reader does not need to interpret a result
The higher-dimensional MINERvA driver (App.~\ref{app:codebase}) extends the 2D driver using an axis list and the same 2D helpers. The first two axes are always $(\pT,\pPar)$, carrying the fiducial phase-space cut and the per-$\pT$ flux. Higher-dimensional runs extend the muon observables $(\pT,\pPar)$ with $\Eavail$, $q_{3}$, $W$ in the phase space of \S\ref{sec:selection}. Every production, universe-sweep, and closure run follows these conventions:
\begin{itemize}[itemsep=2pt]
  \item \textbf{Reco-space masking.} The measured events entering step-1 training are exactly those passing the reconstructed-level selection; nothing outside the reconstructed phase space reaches the classifier.
  \item \textbf{Step-1 training within the phase space.} Step 1 targets the background-subtracted data of \S\ref{sec:method-algo}, purity-weighted or with background simulation injected at negative weight (refined to non-negative weights in the refined configurations). Genuine backgrounds \emph{and} out-of-phase-space signal are removed from the data side. The simulated reconstructed sample contains no out-of-phase-space signal by construction, so neither side of the step-1 training contains events whose truth lies outside the measurement phase space.
        % Accuracy fix (unscored; director ruling 2026-10-26): the list covers every production run, and the 2D result
        % and the 3D and 4D central values use refined negative-weight subtraction (sec:method-algo). Mechanism:
        % 2d-unfolding/unfold_2d_omnifold_unbinned.py:1362-1377 and :1825-1870 @ ddf16f10; nd-unfolding/unfold_nd_omnifold_unbinned.py:619-630 and :900-1030 @ ddf16f10.
        % In the negative-weight modes the signal fakes are injected with the background (2d-unfolding/unfold_2d_omnifold_unbinned.py:1846-1858 @ ddf16f10).
  \item At histogramming, each simulated truth event carries its final OmniFold push-weight times its central-value simulation weight, including MINERvA Tune~v1 reweighting. No efficiency-histogram division is applied: the step-1 regression for unreconstructed truth events already gives an efficiency-corrected truth spectrum, and dividing again would double-correct.
        % Development history: a failure mode that occurred once in
        % development and was caught by the marginal self-validation gate
        % (\S\ref{sec:validation}).
  \item \textbf{Closure mirrors production.} In closure tests the pseudo-data are simulated events passing both the reconstructed and truth selections, weighted identically to the production data path, so closure exercises the same code path as the data unfolding.
\end{itemize}

\paragraph{Step-1 weights for missed truth events.}
% [records reorg] former heading: \subsection{From Part 2, secmethod.tex: the paragraph beginning ``item textbfStep 1 (detector space). A''}
% From sec_method.tex:9 @d55f7aad6; qualifier detail or procedure (DECISION-20261120a, second pass)
% [lead-in, Option B] In the paragraph as the note printed it, this text followed: ``The weight is \emph{pulled back} directly to its truth record.''\par
Truth events without a reconstructed partner (``misses'') receive step-1 weights from a regressor trained on the reconstructed subsample, $w(\mathrm{truth\ features})$. This native OmniFold acceptance treatment requires no separate efficiency correction.

\paragraph{Adding observables in unbinned OmniFold.}
% [records reorg] former heading: \subsection{From Part 2, secmethod.tex: the paragraph beginning ``Unbinned OmniFold adds an observable by''}
% From sec_method.tex:14 @d55f7aad6; qualifier detail or procedure (DECISION-20261120a, second pass)
Unbinned OmniFold adds an observable by adding an input to the step-1 and step-2 classifiers, without an explicit migration matrix or per-bin response tensor. It avoids the combinatorial binning penalty of matrix unfolding but still has a dimensional cost: classifier training, acceptance and support coverage, sample sparsity and uncertainty quantification become more demanding as observables are added. The primary methodological motivation is the absence of a practical analog of D'Agostini iterative Bayesian unfolding (IBU).

\paragraph{Gradient-boosted decision trees as classifiers.}
% [records reorg] former heading: \subsection{From Part 2, secmethod.tex: the paragraph beginning ``Gradient-boosted decision trees (GBDTs) replace OmniFold's''}
% From sec_method.tex:50 @d55f7aad6; qualifier detail or procedure (DECISION-20261120a, second pass)
Gradient-boosted decision trees (GBDTs) replace OmniFold's original neural networks when unfolding scalar observables. With at most five high-level features per event, tree ensembles typically match or outperform neural networks~\cite{Grinsztajn:2022tabular} and train deterministically and far faster on a central processing unit (CPU). All unfolding results here use GBDTs (point-cloud representations, App.~\ref{sec:pet}). Calibration tests and multilayer perceptron (MLP) cross-checks test the implementation (classifier checks, \S\ref{sec:3d-classifier}).

\paragraph{Learner of the quoted 2D central value.}
% [records reorg] former heading: \subsection{From Part 2, secmethod.tex: the paragraph beginning ``The unfolding code supports interchangeable learners''}
% From sec_method.tex:52 @d55f7aad6; qualifier detail or procedure (DECISION-20261120a, second pass)
The unfolding code supports interchangeable learners (component terminology and code paths in App.~\ref{app:codebase}). The quoted 2D central value and all its uncertainty ensembles use LightGBM with refined negative-weight subtraction (background method, App.~\ref{app:negweight}). The estimator matrix (Table~\ref{tab:estimators-matrix}) lists this as the \estTwoD{}. The ensembles cover the systematic-universe sweep, Poisson bootstrap, training-seed variations and run-period unfoldings.
% [records reorg] former heading: \subsection{From Part 2, secmethod.tex: the paragraph beginning ``The unfolding code supports interchangeable learners''}
% From sec_method.tex:52 @d55f7aad6; qualifier detail or procedure (DECISION-20261120a, second pass)
% [lead-in, Option B] In the paragraph as the note printed it, this text followed: ``The step-1 classifier, the step-2 classifier and the regressor for unreconstructed truth events share one setting: $100$ trees, $8$ leaves and learning rate $0.1$.''\par
The universe sweep supplies the central value at seed~42. Thus the systematic covariance shares its estimator and seed (systematic universes, \S\ref{sec:uq-universes}). The quoted statistical covariance also uses seed~42, with three independent bootstrap streams (data, signal and background simulation, \S\ref{sec:uq-boot}). The training-seed covariance uses five seed settings that include seed~42 (\S\ref{sec:uq-ml}).

\paragraph{Reproducibility and seed dispersion of the central value.}
% [records reorg] former heading: \subsection{From Part 2, secmethod.tex: the paragraph beginning ``A central unfolding at seed45 differs''}
% From sec_method.tex:54 @d55f7aad6; qualifier detail or procedure (DECISION-20261120a, second pass)
In the paragraph as the note printed it, this text followed: ``A central unfolding at seed~45 differs from the seed-42 one by at most \SI{1.4}{\percent} in a cell.''\par
The quoted central value has a recorded configuration and repeats numerically. It agrees with the refined 2D data unfolding (background study, App.~\ref{sec:negweight-footing}).
% carried with the text from sec_method.tex:55 @d55f7aad6
% Source: docs/publication/DECISION-20261011-joseph-2d-central-lightgbm-seed42.md; receipts
% [records reorg] former heading: \subsection{From Part 2, secmethod.tex: the paragraph beginning ``A central unfolding at seed45 differs''}
% From sec_method.tex:54 @d55f7aad6; qualifier detail or procedure (DECISION-20261120a, second pass)
% [lead-in, Option B] In the paragraph as the note printed it, this text followed: ``The purity-subtracted seed-42 central value has no recorded run configuration, although an independent rerun and a regeneration whose run configuration is recorded both reproduce it (Ref.~\cite{NoteRecords2026}, \recentry{2d-central-value}).''\par
Against the ten purity-subtracted LightGBM central values at seeds 1--10 it is mildly over-dispersed (root-mean-square standardized difference 1.28, maximum 4.25). The choice of seed~42 is recorded in App.~\ref{app:devhistory}.
% carried with the text from sec_method.tex:55 @d55f7aad6
% Source: docs/publication/DECISION-20261011-joseph-2d-central-lightgbm-seed42.md; receipts

\paragraph{Multidimensional estimator settings.}
% [records reorg] former heading: \subsection{From Part 2, secmethod.tex: the paragraph beginning ``paragraphHyperparameter choice. Neither OmniFold classifier nor''}
% From sec_method.tex:144 @d55f7aad6; qualifier detail or procedure (DECISION-20261120a, second pass)
The \estMultiD{} keeps five iterations under the precision study's rule (\S\ref{sec:gbdt-model-dependence}). Its refinement classifier, 400 trees and 31 leaves, was selected before the joint-test design was fixed (\S\ref{sec:joint5d}).

\paragraph{Fixed GBDT settings and capacity.}
% [records reorg] former heading: \subsection{From Part 2, secmethod.tex: the paragraph beginning ``The fixed settings are the scikit-learn''}
% From sec_method.tex:146 @d55f7aad6; qualifier detail or procedure (DECISION-20261120a, second pass)
The fixed settings are the scikit-learn \texttt{GradientBoosting} defaults: 100 trees, learning rate $0.1$, and depth-3 capacity. The LightGBM learner of the central value and its ensembles uses matched capacity ($\texttt{num\_leaves}=8=2^{3}$), so both have the same deliberately modest capacity, an OmniFold regularizer (algorithm, \S\ref{sec:method-algo}).

\paragraph{Keras multilayer-perceptron cross-check.}
% [records reorg] former heading: \subsection{From Part 2, secmethod.tex: the paragraph beginning ``The Keras cross-check (Srefsec:3d-classifier) uses the''}
% From sec_method.tex:150 @d55f7aad6; qualifier detail or procedure (DECISION-20261120a, second pass)
The Keras cross-check (\S\ref{sec:3d-classifier}) uses the same multilayer perceptron for both classifiers and missed-event regression. It has three hidden layers of 64, 128 and 64 units and Gaussian error linear unit (GELU) activation. Training uses Adam for at most 50 epochs with early stopping, per-fit input standardization and a fixed shuffle. Settings stay fixed. Agreement does not address sensitivity to single-precision (float32) ties (departure tests, \S\ref{sec:val-departure}).

\paragraph{Performance diagnostics.}
% [records reorg] former heading: \subsection{From Part 2, secmethod.tex: the paragraph beginning ``The following diagnostics report performance. They''}
% From sec_method.tex:152 @d55f7aad6; qualifier detail or procedure (DECISION-20261120a, second pass)
% [lead-in, Option B] In the paragraph as the note printed it, this text followed: ``The following diagnostics report performance. They do not validate the quoted estimator or its intervals. They include closure tests (\S\ref{sec:validation}) and classifier calibration (\S\ref{sec:3d-classifier}).''\par
The exact-split cross-check measures insensitivity to doubling the iterations (below). Method-dependence studies report backend agreement (\S\ref{sec:results-tension}). The training-seed variation band measures residual machine-learning (ML) stochasticity (\S\ref{sec:systematics}).

\paragraph{Iteration count.}
% [records reorg] former heading: \subsection{From Part 2, secmethod.tex: the paragraph beginning ``Production uses n=5 iterations. Between five''}
% From sec_method.tex:154 @d55f7aad6; qualifier detail or procedure (DECISION-20261120a, second pass)
Production uses $n=5$ iterations.
% [records reorg] former heading: \subsection{From Part 2, secmethod.tex: the paragraph beginning ``Production uses n=5 iterations. Between five''}
% From sec_method.tex:154 @d55f7aad6; qualifier detail or procedure (DECISION-20261120a, second pass)
% [lead-in, Option B] In the paragraph as the note printed it, this text followed: ``Between five and ten iterations, the exact-split total cross section agrees to within \SI{0.026}{\percent} (validation, \S\ref{sec:validation}). That measures movement of this diagnostic for that cross-check. It does not establish convergence of every functional.''\par
Separately, the unfolded result reproduces the integrated cross section of an independently published binned measurement of the same data (results, \S\ref{sec:results}). That comparison is not an iteration test.

\paragraph{Inputs to the step-1 classifier.}
% [records reorg] former heading: \subsection{From Part 2, secmethod.tex: the paragraph beginning ``Every production, universe-sweep and closure run''}
% From sec_method.tex:164 @d55f7aad6; qualifier detail or procedure (DECISION-20261120a, second pass)
% [lead-in, Option B] In the paragraph as the note printed it, this text followed: ``Every production, universe-sweep and closure run follows the same pipeline conventions (records document, Ref.~\cite{NoteRecords2026}).''\par
Only events passing reconstructed-level selection reach the step-1 classifier.
% [records reorg] former heading: \subsection{From Part 2, secmethod.tex: the paragraph beginning ``Every production, universe-sweep and closure run''}
% From sec_method.tex:164 @d55f7aad6; qualifier detail or procedure (DECISION-20261120a, second pass)
% [lead-in, Option B] In the paragraph as the note printed it, this text followed: ``Both sides of step-1 training exclude events with truth outside the measurement phase space.''\par
Closure pseudo-data follow the production data path.

\paragraph{Class weights and effective sample size.}
% [records reorg] former heading: \subsection{From Part 2, secmethod.tex: the paragraph beginning ``paragraphTraining details. Each classifier is trained''}
% From sec_method.tex:17 @fdd67e926; training detail (DECISION-20261120a records route; the director 2026-11-20)
The two classes keep their physical weights, with no class balancing and no renormalization between iterations. The learned ratio $p/(1-p)$ therefore carries the normalization, with $p$ restricted to $[10^{-6},1-10^{-6}]$.
% [records reorg] former heading: \subsection{From Part 2, secmethod.tex: the paragraph beginning ``paragraphTraining details. Each classifier is trained''}
% From sec_method.tex:17 @fdd67e926; training detail (DECISION-20261120a records route; the director 2026-11-20)
% [lead-in, Option B] In the paragraph as the note printed it, this text followed: ``In the refined configurations, a classifier separating positive from negative weights, trained with $|w|$, gives each event the weight $|w|\max(0,2g-1)$.''\par
For the \estMultiD{}, the per-event Kish effective sample size and weight tails have been computed on the sixteen reported cells~\cite{UnfoldedWeightsESS2026}. They describe the unfolded simulated-event weights and are not an uncertainty, coverage or validity statement.

\paragraph{Check-stage moves at the journal-voice rewrite (2026-11-21).}
% [records reorg] former heading: \subsection{Check-stage moves at the journal-voice rewrite (2026-11-21)}
% Limitation-register rows whose sentence the rewritten note (982945a0a) no longer states, moved verbatim from the
% previous note (draft e89694209) under DECISION-20261120a; one move-manifest row each (LR-K ids in
% history-cuts/move-manifest-lr-check-20261121.tsv). The rewritten note states each result's headline limit.

\paragraph{Alternative-seed central unfolding.}
% [records reorg] former heading: \subsubsection{From Part 2, the training settings (row n07)}
% From sec_method.tex:54 @e89694209; register row, check stage 2026-11-21
% [lead-in, Option B] In the paragraph as the note printed it, this text followed: ``The step-1 classifier, the step-2 classifier and the regressor for unreconstructed truth events share one setting: $100$ trees, $8$ leaves and learning rate $0.1$.''\par
A central unfolding at seed~45 differs from the seed-42 one by at most \SI{1.4}{\percent} in a cell.

\paragraph{Disclosures cut at the journal-voice rewrite (check stage, 2026-11-21).}
% [records reorg] former heading: \subsection{Disclosures cut at the journal-voice rewrite (check stage, 2026-11-21)}
% Qualifications and post-hoc accounts that the check-stage claims reads (Grok, codex2) found the rewritten note no longer
% states, moved verbatim from the previous note (draft e89694209); the rewritten note keeps each result's headline limit.
% One move-manifest row per sentence (LR-K ids, history-cuts/move-manifest-lr-check-20261121.tsv).

\paragraph{Missed-event regressor variants and transfer probe.}
% [records reorg] former heading: \subsubsection{From Part 2, the missed-event regressor and the transfer probe}
% From sec_method.tex:17 @e89694209; check-stage claims read
% absence-ok: scopes what RSN's WR tested (its record's scope: "does not attribute F0b's flags"; the half-sample flags' attribution is T3's, not run), as the paper's limitation 3 at 5bbb9d269 states it
% [lead-in, Option B] In the paragraph as the note printed it, this text followed: ``Truth events without a reconstructed partner receive the prediction of an unweighted regressor of the step-1 weight on the truth observables.''\par
A regressor weighted with each event's current weight instead moves the sixteen-cell data result (\S\ref{sec:fived16}) by at most \fvval{rsn-wr-sig-max} of a cell's reported standard deviation. In noise-free, signal-only closures on the full simulation, it leaves the bias flags at the GiBUU hybrid truth and the $q_3$ tilt unchanged. Its effect on the flags from half-sample pseudo-data with background and noise was not tested. A held-out transfer probe among matched events at fixed truth places no cell beyond half its reported standard deviation, at the data or at those two truths. It cannot show that the missed events resemble matched events in what the five features omit.

\paragraph{Hyperparameter choice.}
% [records reorg] former heading: \subsubsection{From Part 2, the hyperparameter choice}
% From sec_method.tex:144 @e89694209; check-stage claims read
Neither OmniFold classifier nor the missed-event regressor underwent a hyperparameter search on data.

\paragraph{Qualifications from the second check-stage claims read (2026-11-21).}
% [records reorg] former heading: \subsection{Qualifications from the second check-stage claims read (2026-11-21)}
% Moved verbatim from the previous note (draft e89694209); LR-K ids in history-cuts/move-manifest-lr-check-20261121.tsv.

\paragraph{Check-stage moves at the journal-voice rewrite (2026-11-21).}
% [records reorg] former heading: \subsection{Check-stage moves at the journal-voice rewrite (2026-11-21)}
% Limitation-register rows whose sentence the rewritten note (982945a0a) no longer states, moved verbatim from the
% previous note (draft e89694209) under DECISION-20261120a; one move-manifest row each (LR-K ids in
% history-cuts/move-manifest-lr-check-20261121.tsv). The rewritten note states each result's headline limit.

\paragraph{Disclosures cut at the journal-voice rewrite (check stage, 2026-11-21).}
% [records reorg] former heading: \subsection{Disclosures cut at the journal-voice rewrite (check stage, 2026-11-21)}
% Qualifications and post-hoc accounts that the check-stage claims reads (Grok, codex2) found the rewritten note no longer
% states, moved verbatim from the previous note (draft e89694209); the rewritten note keeps each result's headline limit.
% One move-manifest row per sentence (LR-K ids, history-cuts/move-manifest-lr-check-20261121.tsv).

\paragraph{Qualifications from the second check-stage claims read (2026-11-21).}
% [records reorg] former heading: \subsection{Qualifications from the second check-stage claims read (2026-11-21)}
% Moved verbatim from the previous note (draft e89694209); LR-K ids in history-cuts/move-manifest-lr-check-20261121.tsv.

\subsection{Background subtraction in the unfolding}\label{recsec:method-background}

\paragraph{Background subtraction (App.~B): detail.}
% Moved verbatim from app_negweight.tex @ a1e6ec190 in the compression under DECISION-20261120a (qualifier and procedure detail
% may move to the records with a pointer); move-manifest-je JE-20.
% [records reorg] former heading: \subsection{Background subtraction (App.~B): detail moved from the note}

\paragraph{Absence of injected negative weights.}
% [records reorg] former heading: \subsection{From Part 2, secmethod.tex: the paragraph beginning ``paragraphBackground subtraction. In binned purity subtraction,''}
% From sec_method.tex:21 @d55f7aad6; qualifier detail or procedure (DECISION-20261120a, second pass)
No negative-weight events are injected.

\paragraph{Refined and binned subtraction variants.}
% [records reorg] former heading: \subsection{From Part 2, secmethod.tex: the paragraph beginning ``For other experiments' subtraction methods, see''}
% From sec_method.tex:32 @d55f7aad6; qualifier detail or procedure (DECISION-20261120a, second pass)
For other experiments' subtraction methods, see \S\ref{sec:intro}.
% [records reorg] former heading: \subsection{From Part 2, secmethod.tex: the paragraph beginning ``For other experiments' subtraction methods, see''}
% From sec_method.tex:32 @d55f7aad6; qualifier detail or procedure (DECISION-20261120a, second pass)
In the paragraph as the note printed it, this text followed: ``A fully unbinned alternative adds negative-weight background simulation on the measured side.''\par
Its bin-integrated target is $(D_{b}-B_{b})/S_{b}$ (data, background and signal).
% [records reorg] former heading: \subsection{From Part 2, secmethod.tex: the paragraph beginning ``For other experiments' subtraction methods, see''}
% From sec_method.tex:32 @d55f7aad6; qualifier detail or procedure (DECISION-20261120a, second pass)
In the paragraph as the note printed it, this text followed: ``The alternative reproduces the binned correction to \SI{\nwPctTot}{\percent} on the total two-dimensional (2D) cross section and its systematic and statistical covariances to within \SI{2}{\percent} (subtraction comparison, App.~\ref{app:negweight}).''\par
It removes residual binning dependence and suits the extended-fiducial analysis with larger backgrounds. The quoted 2D result and covariance, 3D and 4D central values, sixteen-cell result and covariance, and studies and tests of the \estLowCap{} and \estMultiD{} use the refined variant. All other quoted cross sections and uncertainties, including the \resFiveAxis{} of the \estPurity{} and the extended-fiducial value, use binned subtraction. The matrix, Table~\ref{tab:estimators-matrix} (configurations, \S\ref{sec:method-config}), links results to configurations.

\paragraph{Estimators of the multidimensional results.}
% [records reorg] former heading: \subsection{From Part 2, secmethod.tex: the paragraph beginning ``The quoted 2D cross section and''}
% From sec_method.tex:48 @d55f7aad6; qualifier detail or procedure (DECISION-20261120a, second pass)
% [lead-in, Option B] In the paragraph as the note printed it, this text followed: ``The quoted 2D cross section and covariance use the \estTwoD{}.''\par
The 3D and 4D central values, the \resSixteen{} and the generator tests use the \estMultiD{}. Both subtract background with negative weights and a 400-tree refinement classifier. They differ in training histogram construction in the LightGBM learner and in deterministic settings. An earlier \estPurity{} is retained for the \resFiveAxis{} and for specified figures and closures. Three other configurations serve only as diagnostic controls in the studies listed in Table~\ref{tab:estimators-matrix}, the matrix of estimator uses and implementations.

\subsection{Estimators, configurations and the exact-split cross-check}\label{recsec:method-estimators}

\paragraph{The exact-split cross-check's eigenmode decomposition.}
% Move pair H3-020 (director-helper-3; move-manifest-h3-20261117.tsv sends): sec_results.tex:148 @ 79a934a41, sentences 2-4,
% verbatim with the source comments :153-154; sec:results-tension keeps sentence 1 and points here.
% [records reorg] former heading: \subsubsection{The exact-split cross-check's eigenmode decomposition (from \S\ref{sec:results})}
\label{app:devhistory-results}
This is the exact-split cross-check's part of the eigenmode decomposition of \S\ref{sec:results-tension}. For the exact-split cross-check, $\chi^{2}$ is larger and distributed differently: its components of magnitude $5$--$6$ lie in moderate-$\lambda$ modes on the low-$\pT$ ridge of the cross-section peak. Its 10 smallest-$\lambda$ modes carry \SI{3}{\percent} of the $\chi^{2}$, and the 73 best-measured directions give $\chi^{2}=291$ ($\chi^{2}/205=1.42$, again divided by the number of reported bins; App.~\ref{sec:tension}). Going from 5 to 10 iterations changes the LightGBM (seed~1) value by a negligible $+0.01$.
% Revision D' 2026-11-06 (referee reads of D, both ACCEPT; Sonnet and the interpreter: print the raw chi2 and the directions kept): docs/orchestration/receipts/2d-lgbm42-20261011/tension_report_exact_control.txt:56 prints "1e-04  1.419  73". Arithmetic: 1.419 x 205 = 290.9; 205 x [1.4185, 1.4195) is in [290.8, 291.0], 291 at three figures.
% Director-added wording correction 2026-11-06, for Joseph's review before merge (not among the two he approved): the truncated chi2 was printed as chi2/ndf, but its divisor is the 205 reported bins, not the directions kept. 2d-unfolding/diagnose_tension.py @ c3dbe30f4 (main): :110 n = int(mask.sum()), the 205 reported bins; :183 and :193 divide every truncated chi2 by n (the same divisor at 56e30b9c9 :108/:181/:191 and at 735c909da). Receipt docs/orchestration/receipts/2d-lgbm42-20261011/tension_report_exact_control.txt:56 ("1e-04  1.419  73"). No number changed.
% Moved verbatim (compression K9, class H: a superseded configuration's detail; move-manifest-je JE-10.028..).

\paragraph{Seed dependence of the exact-split result.}
% [records reorg] former heading: \subsection{From Part 2, secmethod.tex: the paragraph beginning ``paragraphThe exact-split cross-check. The exact-split configuration''}
% From sec_method.tex:85 @d55f7aad6; qualifier detail or procedure (DECISION-20261120a, second pass)
Its seed and run log appear in the development history (App.~\ref{app:devhistory}). The result depends on the estimator seed. A reseeded run (\texttt{random\_state} 1/2/3) moves every bin (median \SI{0.18}{\percent}, maximum \SI{1.6}{\percent}) and gives $\chi^{2}/\mathrm{ndf}=\chiPaperExactSeeded$ (chi-square per degree of freedom). The spread over seeds is unmeasured.
% carried with the text from sec_method.tex:86 @d55f7aad6
% Search: exact-split seed scan; spread over seeds; matched covariance; exact-split bootstrap; P09b @ ddf16f10
% Search: exact-backend seed scan; exact-split seed scan; P09b (pairings.tsv: "no exact-backend seed scan exists") @ 1a015eef
% [records reorg] former heading: \subsection{From Part 2, secmethod.tex: the paragraph beginning ``paragraphThe exact-split cross-check. The exact-split configuration''}
% From sec_method.tex:85 @d55f7aad6; qualifier detail or procedure (DECISION-20261120a, second pass)
% [lead-in, Option B] In the paragraph as the note printed it, this text followed: ``The cross-check has no matched covariance.''\par
No exact-split bootstrap, universe sweep or seed scan exists.
% carried with the text from sec_method.tex:86 @d55f7aad6
% Search: exact-split seed scan; spread over seeds; matched covariance; exact-split bootstrap; P09b @ ddf16f10
% Search: exact-backend seed scan; exact-split seed scan; P09b (pairings.tsv: "no exact-backend seed scan exists") @ 1a015eef
% carried with the text from sec_method.tex:86 @d55f7aad6
% Search: exact-split seed scan; spread over seeds; matched covariance; exact-split bootstrap; P09b @ ddf16f10
% Search: exact-backend seed scan; exact-split seed scan; P09b (pairings.tsv: "no exact-backend seed scan exists") @ 1a015eef

\paragraph{Estimator configurations.}
% [records reorg] former heading: \subsection{From Part 2, secmethod.tex: the paragraph beginning ``paragraphThe estimator configurations. The six estimator''}
% From sec_method.tex:102 @d55f7aad6; qualifier detail or procedure (DECISION-20261120a, second pass)
Each result names its estimator.
% [records reorg] former heading: \subsection{From Part 2, secmethod.tex: the paragraph beginning ``paragraphThe estimator configurations. The six estimator''}
% From sec_method.tex:102 @d55f7aad6; qualifier detail or procedure (DECISION-20261120a, second pass)
``Deterministic'' means fixed seeds and reproducible training.
% [records reorg] former heading: \subsection{From Part 2, secmethod.tex: the paragraph beginning ``paragraphThe estimator configurations. The six estimator''}
% From sec_method.tex:102 @d55f7aad6; qualifier detail or procedure (DECISION-20261120a, second pass)
In the paragraph as the note printed it, this text followed: ``Results from different configurations are not combined.''\par
The \estTwoD{} adds refined negative-weight subtraction to the nominal LightGBM configuration above. It uses the \estMultiD{} refinement classifier at the estimator seed plus three. Its quoted seed-42 central value repeats numerically (above). The \estPurity{} is the same learner with purity-based background subtraction (algorithm, \S\ref{sec:method-algo}). The \estMultiD{}, the estimator of the generator tests, is a deterministic LightGBM variant with the unbinned negative-weight background treatment (background method, App.~\ref{app:negweight}). Its refinement classifier has 400 trees and 31 leaves.
% [records reorg] former heading: \subsection{From Part 2, secmethod.tex: the paragraph beginning ``paragraphThe estimator configurations. The six estimator''}
% From sec_method.tex:102 @d55f7aad6; qualifier detail or procedure (DECISION-20261120a, second pass)
In the paragraph as the note printed it, this text followed: ``The negative-weight treatment reduces the purity method's nominal-truth bias in the highest-$W$ (hadronic invariant mass) column by about an order of magnitude (from 3.5--\SI{4.1}{\percent} with purity subtraction to 0.11--\SI{0.24}{\percent} with negative weights).''\par
This reaches the no-background closure level. The larger refinement classifier removes most of the remaining background-related bias (background and departure tests, \S\S\ref{sec:val-negweight}--\ref{sec:val-departure}). The \estMultiD{} was selected before the joint-test design was fixed (\S\ref{sec:joint5d}). The \estExact{} is described above.

\paragraph{Histogram construction and deterministic training.}
% [records reorg] former heading: \subsection{From Part 2, secmethod.tex: the paragraph beginning ``The estPurity and the estTwoD leave''}
% From sec_method.tex:148 @d55f7aad6; qualifier detail or procedure (DECISION-20261120a, second pass)
The \estPurity{} and the \estTwoD{} leave histogram construction at its default. Bin edges come from the library's default row sample. The training seed affects the result through that sample. The configurations this section calls deterministic set those edges from every training row and enable the library's deterministic training.

\paragraph{Regularization scan of the $\chi^{2}$.}
% [records reorg] former heading: \subsection{From App. A, appstatmethods.tex: the paragraph beginning ``For the exact-split cross-check and an''}
% From app_statmethods.tex:440 @d55f7aad6; qualifier detail or procedure (DECISION-20261120a, second pass)
Each method spans roughly two decades of $\chi^{2}/205$ across a two-decade scan of its regularization parameter, with no defensible single cutoff. The useful conclusion is diagnostic, not a number: the naive inverse is dominated by tiny, poorly-determined modes, and in the retained high-variance modes the central-value difference is not large. We therefore report direct per-bin pulls (\S\ref{sec:pull}) and this scan rather than a single unqualified $\chi^{2}_{\ours}$.

\paragraph{Separation of configurations.}
% [records reorg] former heading: \subsubsection{From Part 2, the estimator configurations (row n09)}
% From sec_method.tex:102 @e89694209; register row, check stage 2026-11-21
% [lead-in, Option B] In the paragraph as the note printed it, this text followed: ``All use five OmniFold iterations and gradient-boosted trees of nominal capacity (100 trees, 8 leaves) in both OmniFold steps.''\par
Results from different configurations are not combined.

\paragraph{Scope of reduced prior dependence.}
% [records reorg] former heading: \subsubsection{From Part 2, related developments and estimator scope (row n13)}
% From sec_method.tex:201 @e89694209; register row, check stage 2026-11-21
% absence-ok: verbatim move of a register sentence (a scoping statement, not an absence of a record)
% [lead-in, Option B] In the paragraph as the note printed it, this text followed: ``See the records for the distinction from Profile OmniFold (Ref.~\cite{NoteRecords2026}).''\par
Reduced simulation-prior dependence does not establish detector-model independence, adequate phase-space support, or coverage for this result.

\paragraph{Exact-split central-value differences.}
% [records reorg] former heading: \subsubsection{From Part 2, the exact-split cross-check}
% From sec_method.tex:85 @e89694209; check-stage claims read
% [lead-in, Option B] In the paragraph as the note printed it, this text followed: ``The exact-split configuration (\texttt{exact}) uses scikit-learn \texttt{GradientBoosting} with single-threaded exact-split classification and regression trees (CART).''\par
The exact-split central value differs from the purity-subtracted LightGBM seed-42 central value by a per-bin median \SI{0.98}{\percent} (maximum \SI{12.2}{\percent}), 1.32 statistical standard deviations. It differs from the quoted central value, which also differs from it in background treatment, by a per-bin median of \SI{1.03}{\percent} (maximum \SI{17.6}{\percent}).
% [records reorg] former heading: \subsubsection{From Part 2, the scope of exact-split studies}
% From sec_method.tex:103 @e89694209; check-stage claims read
In the paragraph as the note printed it, this text followed: ``Studies labeled \emph{exact-split} use that cross-check, an earlier quoted 2D central value.''\par
They have not been repeated for the quoted estimator.

\paragraph{Shape distance of the exact-split cross-check.}
% [records reorg] former heading: \subsubsection{From App. A, the exact-split shape distance}
% From app_statmethods.tex:398 @e89694209; check-stage claims read
For the exact-split cross-check, which as above is not recomputed for the LightGBM central value, with the paper's \texttt{TotalCovariance} (Jacobian-propagated) the shape $\chi^{2}/\mathrm{ndf}$ on $205$ bins is $\boxed{3.596}$: essentially the absolute-normalization $3.661$ \emph{minus} the contribution of overall scale. The total-cross-section ratio $\sigma_{\tot}^{\ours}/\sigma_{\tot}^{\paper} = 1.0111$ is a $1\%$ scale difference. The remaining residual comes from $\sim 1\%$ shape disagreement (median bin ratio $1.0064$; $94\%$ of bins within $10\%$ of the paper).

\paragraph{Relative uncertainties with respect to the exact-split cross-check.}
% [records reorg] former heading: \subsection{Relative uncertainties with respect to the exact-split cross-check (from Part 3, the uncertainty budget)}
% From sec_results.tex:30 @ draft 47c318849, the paragraph's last sentence, with its source comment
Taken with respect to the exact-split cross-check instead, the same combined covariance has a per-bin median of \SI{6.468}{\percent}. % G2R: RE-2-shapes run 2 (C_new), g2r-integration/ledger.tsv g2r-res-q-vs-ec; a source comment citing RE-2 part C for this value gives the earlier value

\subsection{Binning, migrations and the reported product}\label{recsec:method-binning}

\paragraph{Structure of the truth-to-reco migrations.}
% [records reorg] former heading: \subsection{Structure of the truth-to-reco migrations (from Part 2, sec:method-binning)}
% From sec_method.tex (dfc41cdf5) (draft f6a4fc372 lineage); description of the migration figure beyond its caption
The structure differs qualitatively between the muon and hadron sides. The spectrometer muon measurement makes the $\pT$ and $\pPar$ maps nearly tridiagonal, with migrations rarely exceeding one bin. Reconstruction stability dips to ${\sim}0.45$--$0.5$ for $\pT$ and ${\sim}0.35$ for $\pPar$ where bins are finest, recovering in wide edge bins. The calorimetric axes have coarser bins but relatively broader, visibly asymmetric migrations. $\Eavail$ preferentially migrates \emph{downward} because energy carried by neutrons and other unresponsive particles is not collected. Reconstruction stability falls from $0.61$ in the lowest bins to $0.43$ in the highest. $q_{3}$, reconstructed from the muon $Q^{2}$ and the calorimetric recoil, inherits a milder version of the same pattern. $W$ compounds $\Eavail$ and muon smearing and has the broadest migrations; two-bin shifts are common in the resonance region. The lowest-$W$ (quasi-elastic) bin has the highest stability at $0.71$.

\paragraph{No underflow or overflow in the reported product.}
% [records reorg] former heading: \subsection{No underflow or overflow in the reported product (from Part 2, sec:method-binning)}
% From sec_method.tex @1ad08b466; argument a reader does not need to interpret a result
Neither underflow nor overflow appears in the reported product, by construction rather than by truncation: the unbinned unfolding has no response-matrix under- or overflow, and the reporting bins are applied once, to events that all lie inside the grids. Truth-level muon edges coincide with the measurement phase space (\S\ref{sec:selection}); no signal is defined beyond them. Hadronic axes span $[0,100]\,\si{GeV}$, with catch-bin edges far beyond the kinematic reach at $\langle E_\nu\rangle\sim\SI{6}{GeV}$, so the truth grids are exhaustive. A selected reco event outside the grid is treated as unaccepted and excluded from the measured sample by \S\ref{sec:method-driver}'s masking convention. Its in-phase-space truth content is recovered, as for any unreconstructed event, by the treatment of truth events without a reconstructed partner.

\paragraph{Correction factor, self-test and flux integral.}
% [records reorg] former heading: \subsection{From Part 2, secmethod.tex: the paragraph beginning ``In n dimensions, the differential cross''}
% From sec_method.tex:176 @d55f7aad6; qualifier detail or procedure (DECISION-20261120a, second pass)
% [lead-in, Option B] In the paragraph as the note printed it, this text followed: ``Thus the defined truth sample enters and $c\equiv1$ per bin.''\par
Negative truth $W^{2}$ is excluded by the five-observable sample definition (\S\ref{sec:eavailw-definitions}). This differs from losing a reconstructed event from the input. The factor remains a monitored self-check. Dropping unmatched truth events would give $c<1$.
% [records reorg] former heading: \subsection{From Part 2, secmethod.tex: the paragraph beginning ``In n dimensions, the differential cross''}
% From sec_method.tex:176 @d55f7aad6; qualifier detail or procedure (DECISION-20261120a, second pass)
% [lead-in, Option B] In the paragraph as the note printed it, this text followed: ``$\Phi$ is the flux integral by transverse muon momentum $\pT$. It is tabulated per square meter and broadcast over other axes.''\par
The formula uses that value.
% [records reorg] former heading: \subsection{From Part 2, secmethod.tex: the paragraph beginning ``In n dimensions, the differential cross''}
% From sec_method.tex:176 @d55f7aad6; qualifier detail or procedure (DECISION-20261120a, second pass)
% [lead-in, Option B] In the paragraph as the note printed it, this text followed: ``The factor $10^{4}$ converts flux tabulated per square meter to the reported $\si{cm^2}$ ($\SI{1}{m^2}=10^{4}\,\si{cm^2}$), giving units of $\si{cm^2}$ per axis unit per nucleon.''\par
The implementation (code, App.~\ref{app:codebase}) is self-tested to reproduce stored lower-dimensional reference extractions to $<10^{-12}$. Since $\sigma\propto1/\Phi$, $\Phi$ varies per flux universe when building the systematic covariance. This variation dominates the Flux band (systematics, \S\ref{sec:systematics}).

\paragraph{Nucleon count and flux units.}
% [records reorg] former heading: \subsection{From Part 2, secmethod.tex: the paragraph beginning ``For higher-dimensional extractions the same integral''}
% From sec_method.tex:179 @d55f7aad6; qualifier detail or procedure (DECISION-20261120a, second pass)
% [lead-in, Option B] In the paragraph as the note printed it, this text followed: ``For higher-dimensional extractions the same integral is $\SI{6.243e-8}{cm^{-2}}$ per proton on target, from 0 to \SI{100}{GeV} after the neutrino--electron-scattering constraint.''\par
$N$ is the formula's fiducial nucleon count. The $\si{cm^{-2}}$ value is the converted integral. It is not substituted for $\Phi$.

\paragraph{Extension-axis bin edges.}
% [records reorg] former heading: \subsection{From Part 2, secmethod.tex: the paragraph beginning ``The 2D grid has 14 pT''}
% From sec_method.tex:184 @d55f7aad6; qualifier detail or procedure (DECISION-20261120a, second pass)
% [lead-in, Option B] In the paragraph as the note printed it, this text followed: ``The 19 unreported diagonal cells have $\pT^{\mathrm{hi}}/\pPar^{\mathrm{lo}}>\tan20^{\circ}$.''\par
The extension axes use $\Eavail$: $[0,0.1,0.2,0.4,0.8,1.5,3.0,100]\,\si{GeV}$; $q_{3}$ (three-momentum transfer): $[0,0.2,0.4,0.6,0.8,1.2,2.0,100]\,\si{GeV}$; $W$: $[0,1.1,1.4,1.8,2.2,3.0,100]\,\si{GeV}$.
% [records reorg] former heading: \subsection{From Part 2, secmethod.tex: the paragraph beginning ``The 2D grid has 14 pT''}
% From sec_method.tex:184 @d55f7aad6; qualifier detail or procedure (DECISION-20261120a, second pass)
% [lead-in, Option B] In the paragraph as the note printed it, this text followed: ``Wide final bins capture full charged-current (CC) inclusive tails in marginals.''\par
Truncating the top edge would break the marginalization checks in \S\ref{sec:3d-anchor}. Post-unfolding binning lets us revisit these choices without rerunning unfolding.

\paragraph{Truth-to-reco migration.}
% [records reorg] former heading: \subsection{From Part 2, secmethod.tex: the paragraph beginning ``Truthtoreco migration on the analysis binning''}
% From sec_method.tex:187 @d55f7aad6; qualifier detail or procedure (DECISION-20261120a, second pass)
Truth$\to$reco migration on the analysis binning (Fig.~\ref{fig:migration}) shows the smearing OmniFold must undo. OmniFold uses no response matrix. Median reconstruction stability (same-bin fraction, Fig.~\ref{fig:migration}) is $\sim0.6$ on each axis. Such migrations would badly bias a bin-by-bin correction. Unfolding is required.

\subsection{Related developments: Profile OmniFold, SPINUP and gradient-boosted trees}\label{recsec:method-related}

\paragraph{Profile OmniFold~\cite{Zhu:2025profile} and related developments.}
% [records reorg] former heading: \subsection{Profile OmniFold~\cite{Zhu:2025profile} and related developments (from Part 2)}
% From sec_method.tex @1ad08b466; literature context
Its distinct algorithmic construction complements earlier unbinned profiled unfolding~\cite{Chan:2023tbf}.

\paragraph{Gradient-boosted trees in unbinned unfolding.}
% [records reorg] former heading: \subsection{Gradient-boosted trees in unbinned unfolding (from Part 2, sec:method-config)}
% From sec_method.tex @1ad08b466; literature context
Ref.~\cite{Milton:2025tools} demonstrates their effectiveness for unbinned unfolding with few high-level features, complementing deep models for low-level event representations.

\paragraph{Profile OmniFold.}
% [records reorg] former heading: \subsection{From Part 2, secmethod.tex: the paragraph beginning ``Profile OmniFoldciteZhu:2025profile extends classifier-based expectation maximization''}
% From sec_method.tex:199 @d55f7aad6; qualifier detail or procedure (DECISION-20261120a, second pass)
Profile OmniFold~\cite{Zhu:2025profile} extends classifier-based expectation maximization (EM) to forward-model nuisance parameters. Profiling nuisances with observed data would change our inference model, which propagates an external nuisance ensemble through a fixed procedure. Neither a profiled estimator nor its uncertainty enters here.

\paragraph{SPINUP and AUSSIE.}
% [records reorg] former heading: \subsection{From Part 2, secmethod.tex: the paragraph beginning ``SPINUP learns a forward mappingciteButter:2025spinup. AUSSIE''}
% From sec_method.tex:201 @d55f7aad6; qualifier detail or procedure (DECISION-20261120a, second pass)
SPINUP learns a forward mapping~\cite{Butter:2025spinup}. AUSSIE replaces OmniFold's second step with a different loss and avoids the outer iteration~\cite{Ore:2026aussie}. Both offer possible future cross-checks with different algorithmic assumptions.
% [records reorg] former heading: \subsection{From Part 2, secmethod.tex: the paragraph beginning ``SPINUP learns a forward mappingciteButter:2025spinup. AUSSIE''}
% From sec_method.tex:201 @d55f7aad6; qualifier detail or procedure (DECISION-20261120a, second pass)
In the paragraph as the note printed it, this text followed: ``Reduced simulation-prior dependence does not establish detector-model independence, adequate phase-space support, or coverage for this result.''\par
Comparisons would require matched signal, selection, normalization, and uncertainty definitions.
